\documentclass[11pt,a4paper]{article}
\usepackage[T1]{fontenc}
\usepackage[utf8]{inputenc}
\usepackage{lmodern}
\usepackage{amsmath,amssymb}
\usepackage[margin=25mm]{geometry}
\usepackage[hidelinks]{hyperref}
\hypersetup{pdftitle={From quantum time slices to relativistic and Newtonian motion},pdfauthor={navstokgap research working notes}}
\newcommand{\tightlist}{\setlength{\itemsep}{0pt}\setlength{\parskip}{0pt}}
\setlength{\emergencystretch}{3em}
\setcounter{secnumdepth}{0}
\begin{document}
\hypertarget{from-quantum-time-slices-to-relativistic-and-newtonian-motion}{%
\section{From quantum time slices to relativistic and Newtonian
motion}\label{from-quantum-time-slices-to-relativistic-and-newtonian-motion}}

\textbf{Result, 2026-10-27.} The finite-slice construction of
\href{quantum-blocks-construct-newton.md}{Newton motion} extends to a
nonquadratic kinetic energy \(T(P)\) whenever \(T'\) and \(V'\) are
globally Lipschitz. The sharper mean rate stated below also uses the
bounded third derivatives in §1. For the positive-energy relativistic
dispersion \[T_c(p)=\sqrt{m^2c^4+c^2p^2}-mc^2,\qquad m,c>0,\tag{1}\] the
constants in the packet estimates can be chosen independently of
\(c\ge c_0>0\). For fixed initial data and finite time, the
finite-product coherent-state means obey \[
 \sup_{0\le t\le T}\|z_{\hbar,\pi,c}(t)-z_{\rm N}(t)\|_\gamma
 \le C_T\left(\hbar+|\pi|+\frac1{c^2}\right).\tag{2}
\] Here \(z_{\rm N}\) is the Newton solution for the same \(V\), \(\pi\)
is an arbitrary positive time partition and \(\gamma>0\) fixes a
phase-space norm. This is a mean bound; the packet widths remain
\(O(\sqrt\hbar)\). Thus \(\hbar\downarrow0\), time refinement and
\(c\to\infty\) may be taken jointly at the level of these means.
Explicit constants and a proof are given below.

The exact Hamiltonians are self-adjoint and, at each fixed \(\hbar\),
their nonrelativistic limit is strong resolvent, hence strong unitary,
convergence. A separate spectral-Dirichlet calculation gives uniform
position convergence through a free hard-wall bounce, even when
\(\hbar\downarrow0\) and \(c\to\infty\) together. These are
one-particle, positive-energy constructions with an external potential;
they do not construct a Lorentz-covariant interacting field theory or
supply a positive universal action constant.

The coherent-state classical limit and relativistic Feynman--Kac
representations have established precedents; this note assembles a
specific joint bound and makes no claim of priority. In particular,
\href{https://doi.org/10.1016/0022-1236(90)90049-Q}{Carmona--Masters--Simon}
(abstract) treat relativistic Schrödinger operators through processes
with stationary independent increments, and
\href{https://arxiv.org/abs/1109.5829}{Hiroshima--Ichinose--Lőrinczi}
(abstract) give a subordinate-process Feynman--Kac formula in a more
general spin and magnetic setting.

\textbf{Targeted advisor check (GPT-6.1 Sol, 2026-10-27).} The written
packet, domain, strong-resolvent and spectral-wall arguments were
checked with corrections applied. This is not independent peer review of
the result or of the imported operator and probability theorems. The
centered \(O(\hbar)\) improvement is a subsequent written derivation,
not part of that earlier check.

\hypertarget{a-general-dispersion-and-its-quantum-realization}{%
\subsection{1. A general dispersion and its quantum
realization}\label{a-general-dispersion-and-its-quantum-realization}}

Work on \(L^2(\mathbb R)\) with \(Qx=x\), \(P=-i\hbar\,d/dx\), and
\([Q,P]=i\hbar\). Assume for the full operator argument that real
\(T,V\in C^\infty(\mathbb R)\) are bounded below and that all
derivatives of order at least two are bounded. Add constants so that
\(T,V\ge0\), which affects only a phase. Write
\[L=\|V''\|_\infty,\qquad M=\|T''\|_\infty.\tag{3}\] For the sharper
mean estimate also write \(L_3=\|V'''\|_\infty\) and
\(M_3=\|T'''\|_\infty\), which are finite under the standing smoothness
assumptions. The finite-product estimates alone need only \(C^1\)
functions with globally Lipschitz derivatives; the smoother class makes
the domain argument explicit. Define on Schwartz space
\[H_\hbar=T(P)+V(Q),\qquad U_\hbar(t)=e^{-itH_\hbar/\hbar}.\tag{4}\]

To specify the domain, fix units for \(Q,P\) and put \(N=1+Q^2+P^2\).
Since \(T,V\) grow at most quadratically,
\(\|H_\hbar f\|\le C_\hbar\|Nf\|\) on \(\mathcal S\). Moreover \[
 [H_\hbar,N]
 =-i\hbar\{T'(P)Q+QT'(P)\}
  +i\hbar\{V'(Q)P+PV'(Q)\},\tag{5}
\] and the linear growth of \(T',V'\) and boundedness of their
derivatives give \(\|[H_\hbar,N]f\|\le C_\hbar\|Nf\|\). Therefore \[
 |\langle H_\hbar f,N^2f\rangle
      -\langle N^2f,H_\hbar f\rangle|
 \le2C_\hbar\langle f,N^2f\rangle.\tag{6}
\] Schwartz space is a core for \(N^2\). Nelson's commutator theorem
gives essential self-adjointness; the \(N^2\) form estimate also makes
the unitary group leave \(D(N)\) invariant and act strongly continuously
in its graph norm. For the cited energy-limited theorem use
\(G=N^2-\inf\sigma(N^2)\ge0\) and \(E_0=\inf\sigma(N^2)\); the relevant
form domain remains \(D(N)\) and (6) bounds \(G+E_0\). These are
standard imported operator results: see
\href{https://link.springer.com/article/10.1007/s00220-025-05282-w}{Energy-Limited
Quantum Dynamics, Theorems 3.7--3.8 and Corollary B.2} (statements and
hypotheses read). No nonlocal analogue of the local integration-by-parts
graph identity from the earlier proof is assumed.

For \(s\ge0\) put \[K_s=e^{-isV(Q)/\hbar},\quad D_s=e^{-isT(P)/\hbar},
 \quad S_s=D_sK_s.\tag{7}\] On \(D(N)\),
\((S_s-I)f/s\to-iH_\hbar f/\hbar\) and the difference quotients are
bounded from the \(N\) graph domain to \(L^2\). The exact orbit of
\(f\in D(N)\) is compact in that graph norm on \([0,T]\); a telescoping
sum therefore proves that products of \(S_s\) over arbitrary positive
partitions converge strongly to \(U_\hbar(t)\) as the mesh vanishes.
Density and unitarity extend the result to every \(L^2\) state. This is
a fixed-\(\hbar\) assertion, not uniform operator-norm control as
\(\hbar\downarrow0\) or \(c\to\infty\).

\hypertarget{the-finite-slice-spread-and-mean-bounds}{%
\subsection{2. The finite-slice spread and mean
bounds}\label{the-finite-slice-spread-and-mean-bounds}}

Choose \(\gamma>0\), \(h_0>0\), \(0<\hbar\le h_0\), and a coherent
Gaussian with means \((q_0,p_0)\) and initial standard deviations
\[\operatorname{sd}Q=\sqrt{\gamma\hbar/2},\qquad
 \operatorname{sd}P=\sqrt{\hbar/(2\gamma)}.\tag{8}\] These are the usual
coherent widths for \(\gamma=1/(m\omega)\). Use \[
 \|(q,p)\|_\gamma=|q|+\gamma|p|,\quad
 C=\sqrt{2\gamma},\quad
 \beta=M/\gamma+\gamma L,\quad
 a=\max(M/\gamma,\gamma L).\tag{9}
\] When \(a=0\), interpret \((e^{aT}-1)/a\) below as \(T\).

For the state after each complete cell define
\(W_j=\operatorname{sd}Q_j+\gamma\operatorname{sd}P_j\). The exact
finite-factor identities on Schwartz space are
\[K_s^*PK_s=P-sV'(Q),\qquad
 D_s^*QD_s=Q+sT'(P).\tag{10}\] Since the standard deviation of a
Lipschitz function is at most its Lipschitz constant times the original
standard deviation, \[
 W_{j+1}\le(1+sM/\gamma)(1+s\gamma L)W_j,\qquad
 W_j\le C\sqrt\hbar\,e^{\beta T}.\tag{11}
\] The same uniform bound covers the intermediate post-kick state. The
state need not remain Gaussian.

Let \(q_j,p_j\) be the packet means. The exact update reads \[
 p_{j+1}=p_j-s\langle V'(Q)\rangle_j,\qquad
 q_{j+1}=q_j+s\langle T'(P)\rangle_{j,\mathrm{kick}}.\tag{12}
\] Introduce constants independent of the partition and of
\(0<\hbar\le h_0\): \[
 \begin{gathered}
 S=C\sqrt{h_0}\,e^{\beta T},\qquad
 \bar f=|V'(0)|+LS,\qquad
 \bar u=|T'(0)|+(M/\gamma)S,\\
 B=\gamma\bar f+\bar u,\qquad
 R=e^{\beta T}(\|(q_0,p_0)\|_\gamma+TB),\\
 F_T=|V'(0)|+L(R+S),\qquad
 U_T=|T'(0)|+(M/\gamma)(R+S).
 \end{gathered}\tag{13}
\] Equation (12) and (11) give \[
 \|(q_{j+1},p_{j+1})\|_\gamma
 \le e^{\beta s}\|(q_j,p_j)\|_\gamma+sBe^{\beta s};
 \tag{14}
\] hence all vertex means have norm at most \(R\),
\(|\langle V'(Q)\rangle_j|\le F_T\), and
\(|\langle T'(P)\rangle_{j,\mathrm{kick}}|\le U_T\).

Linearly interpolate the means and define
\(\mathcal F_T(q,p)=(T'(p),-V'(q))\). Its Lipschitz constant in
\(\|\cdot\|_\gamma\) is at most \(a\). On every cell, \[
 \|\dot z_{\hbar,\pi}-\mathcal F_T(z_{\hbar,\pi})\|_\gamma
 \le C_\delta|\pi|+C_h\hbar,\qquad
 C_\delta=MF_T+\gamma LU_T,\quad
 C_h=\frac{C^2e^{2\beta T}}2
             \left(\frac{M_3}{\gamma^2}+\gamma L_3\right).\tag{15}
\] For each mean, expand \(V'\) and \(T'\) around its corresponding
current mean; the linear terms have zero expectation. Thus the force
bias is at most \(\frac12L_3\operatorname{Var}_j(Q)\), and the post-kick
velocity bias is at most
\(\frac12M_3\operatorname{Var}_{j,\rm kick}(P)\). Use (11), including
the intermediate-state bound, and add the within-cell changes
\(M|p_{j+1}-p(t)|\) and \(L|q_j-q(t)|\). This proves (15) without
assuming the propagated state is Gaussian.

Comparing two meshes and two positive values \(\hbar,k\), integration
and Gronwall give \[
 \sup_{t\le T}\|z_{\hbar,\pi}(t)-z_{k,\sigma}(t)\|_\gamma
 \le G_T\{C_\delta(|\pi|+|\sigma|)
           +C_h(\hbar+k)\},\quad
 G_T=\frac{e^{aT}-1}{a}.\tag{16}
\] The interpolants are therefore Cauchy as both scales tend to zero,
and their limit solves uniquely
\[\dot q=T'(p),\qquad \dot p=-V'(q).\tag{17}\] No solution is assumed in
defining that limit. At fixed \(\hbar\) the strong product limit plus
the uniform second-moment bounds recovers the exact unitary first
moments: bounded spectral truncations pass by strong convergence and the
tails are bounded by the second moments. Thus (15)--(16) also control
the exact quantum means after the mesh is removed.

\hypertarget{relativistic-dispersion-and-the-newton-limit}{%
\subsection{3. Relativistic dispersion and the Newton
limit}\label{relativistic-dispersion-and-the-newton-limit}}

For (1), \[
 T_c'(p)=\frac{p}{m\sqrt{1+p^2/(m^2c^2)}},\qquad
 0<T_c''(p)=\frac1m
              \left(1+\frac{p^2}{m^2c^2}\right)^{-3/2}
             \le\frac1m.\tag{18}
\] All hypotheses above hold for every \(c>0\), with common \(M=1/m\)
and \(T_c'(0)=0\). For \(c\ge c_0>0\), \[
 |T_c'''(p)|=\frac{3|p|}{m^3c^2}
       \left(1+\frac{p^2}{m^2c^2}\right)^{-5/2}
 \le\frac3{m^2c_0}.\tag{18a}
\] Consequently all constants in (11)--(16) can be chosen independently
of \(c\ge c_0\). The finite-\(c\) classical solution \(z_c\) from (17)
is the positive-energy relativistic particle with
\[\dot q=\frac{pc^2}{\sqrt{m^2c^4+c^2p^2}},\qquad
 \dot p=-V'(q),\qquad |\dot q|<c.\tag{19}\] The external potential is
specified in a chosen frame; (19) does not assert full interacting
Poincaré covariance.

The elementary inequality \(1-(1+x)^{-1/2}\le x/2\) for \(x\ge0\) gives
\[|T_c'(p)-p/m|\le\frac{|p|^3}{2m^3c^2}.\tag{20}\] The Newton trajectory
is already constructed by (17) with \(T_\infty(p)=p^2/(2m)\). The common
bound (13) gives \(|p_{\rm N}(t)|\le R/\gamma\). Comparing its vector
field with (19) and applying Gronwall, \[
 \sup_{t\le T}\|z_c(t)-z_{\rm N}(t)\|_\gamma
 \le G_T\frac{(R/\gamma)^3}{2m^3c^2}.\tag{21}
\] Together with (15)--(16), \[
 \sup_{t\le T}\|z_{\hbar,\pi,c}(t)-z_{\rm N}(t)\|_\gamma
 \le G_T\!\left[
 C_\delta|\pi|+C_h\hbar
 +\frac{(R/\gamma)^3}{2m^3c^2}\right].\tag{22}
\] This proves (2). The bound is for fixed initial phase-space data and
finite \(T\); it is not uniform in momenta of order \(mc\).

For any compactly supported smooth phase-space symbol \(b\), the
centered second-order Weyl Fourier argument in the
\href{groupoid-quantum-newton-states.md}{groupoid companion} converts
the spread and mean bounds into \[
 |\langle\operatorname{Op}_\hbar^W(b)\rangle_j-b(z_{\rm N}(t_j))|
 \le J_{b,\gamma}C^2e^{2\beta T}\hbar+
 I_{b,\gamma}G_T\!\left(
 C_\delta|\pi|+C_h\hbar
 +\frac{(R/\gamma)^3}{2m^3c^2}\right),\tag{23}
\] where \[I_{b,\gamma}=(2\pi)^{-2}\int|\widehat b(\xi,\eta)|
            (|\xi|+|\eta|/\gamma)\,d\xi d\eta,\quad
 J_{b,\gamma}=\frac1{2(2\pi)^2}\int|\widehat b(\xi,\eta)|
            (|\xi|+|\eta|/\gamma)^2\,d\xi d\eta.\tag{24}\] Thus a chosen
diagonal schedule \(\hbar=S_*\epsilon\), \(|\pi|\le T_*\epsilon^\alpha\)
and \(c\ge c_*\epsilon^{-\kappa}\), with \(\alpha,\kappa>0\), gives
\(O(\epsilon+\epsilon^\alpha+\epsilon^{2\kappa})\) for those observables
at partition vertices. The groupoid deformation, mesh and \(c\) remain
distinct until this schedule is chosen.

\hypertarget{fixed-quantum-action-and-the-nonrelativistic-limit}{%
\subsection{4. Fixed quantum action and the nonrelativistic
limit}\label{fixed-quantum-action-and-the-nonrelativistic-limit}}

For fixed \(\hbar>0\), define the nonnegative closed forms \[
 q_c[\psi]=\|T_c(P)^{1/2}\psi\|^2
            +\|V(Q)^{1/2}\psi\|^2,\qquad
 D(q_c)=D(T_c(P)^{1/2})\cap D(V(Q)^{1/2}).\tag{25}
\] The functional calculus expression
\[T_c(p)=\frac{p^2/m}{1+\sqrt{1+p^2/(m^2c^2)}}\tag{26}\] shows that
\(T_c(p)\) increases pointwise to \(p^2/(2m)\) as \(c\to\infty\). The
increasing limit of the forms has domain \(D(P)\cap D(\sqrt V)\) and
equals the nonrelativistic form. The monotone convergence theorem for
closed forms therefore yields
\[H_{\hbar,c}\xrightarrow[c\to\infty]{\rm strong\ resolvent}
 H_{\hbar,\infty}=P^2/(2m)+V(Q).\tag{27}\] See
\href{https://link.springer.com/article/10.1007/s13373-018-0118-0}{Simon,
``Tosio Kato's work,'' Theorem 10.12} (theorem statement read). Standard
spectral convergence then gives strong convergence of the corresponding
unitary groups, uniformly for \(t\) in compact intervals, for each fixed
\(\hbar\). This is not an operator-norm or uniform-in-\(\hbar\) claim.

In Euclidean time the relativistic kinetic semigroup is positive, but
its reference process is no longer Brownian in physical time. Indeed,
for \(a,b>0\) the subordination identity is
\[e^{-a\sqrt b}=\frac a{2\sqrt\pi}\int_0^\infty
 s^{-3/2}e^{-a^2/(4s)}e^{-bs}\,ds.\tag{28}\] With \(a=tc/\hbar\) and
\(b=p^2+m^2c^2\), multiply the right side of (28) by \(e^{tmc^2/\hbar}\)
to obtain \(e^{-tT_c(P)/\hbar}\) as a positive mixture of Gaussian heat
operators. Setting \(p=0\) shows that the mixing measure has unit mass.
The associated subordinate-Brownian process has jumps. Feynman--Kac with
\(V\) is an additional standard construction, not a claim that the
Gaussian midpoint variance or its large-field bound from the preceding
note remains unchanged. For a general \(T\) with bounded \(T''\),
positivity of the Euclidean semigroup need not hold.

\hypertarget{a-relativistic-spectral-wall}{%
\subsection{5. A relativistic spectral
wall}\label{a-relativistic-spectral-wall}}

The hard-wall model has to be specified again because \(T_c(P)\) is
nonlocal. Let \(\Delta_D\) be the Dirichlet Laplacian on \((0,\infty)\)
and define the \textbf{spectral} operator \[H_{\hbar,c}^{D}
 =T_c\!\left(\sqrt{-\hbar^2\Delta_D}\right).\tag{29}\] It is
self-adjoint. Under odd extension it is unitarily equivalent to
\(T_c(P)\) restricted to the odd subspace, since \(T_c\) is even.

Use the odd coherent state of \href{newton-dimock-balaban-guide.md}{the
didactic guide, equation (20)}. Write its overlap as \(\gamma_\hbar\),
its initial full-line widths as \(\sigma_q,\sigma_p\), and put
\[a_c(t)=q_0+tT_c'(p_0),\qquad q_0>0,\quad p_0<0.\tag{30}\] The
full-line Heisenberg identity is \(Q(t)=Q+tT_c'(P)\). For the evolved
Gaussian \(g_{\hbar,t}\), \[
 \|(Q-a_c(t))g_{\hbar,t}\|
 \le\sigma_q+|t|\,\|[T_c'(P)-T_c'(p_0)]g_\hbar\|
 \le\sigma_q+\frac{|t|}{m}\sigma_p.\tag{31}
\] No preservation of Gaussian shape is needed. Apply
\(||x|-|a_c||\le|x-a_c|\) to each of the two terms in the odd state. The
half-line position expectation then satisfies \[
 \left|\langle Q\rangle_{\hbar,c}^{D}(t)-|a_c(t)|\right|
 \le\sqrt{\frac2{1-\gamma_\hbar}}
             \left(\sigma_q+\frac{|t|}{m}\sigma_p\right).\tag{32}
\] Equation (20) also gives, uniformly for \(|t|\le T\), \[
 \left|\langle Q\rangle_{\hbar,c}^{D}(t)
          -\left|q_0+\frac{p_0t}{m}\right|\right|
 \le\sqrt{\frac2{1-\gamma_\hbar}}
             \left(\sigma_q+\frac{T}{m}\sigma_p\right)
      +\frac{T|p_0|^3}{2m^3c^2}.\tag{33}
\] The two classical impact times generally differ, but (33) is uniform
through both. At fixed \(c\), (32) gives the corresponding classical
relativistic elastic bounce; at fixed \(q_0,p_0,T\), (33) is
\(O(\sqrt\hbar+c^{-2})\) as the two limits are taken in any order.

The spectral operator in (29) subordinates the \textbf{killed Brownian}
semigroup. It is distinct from first subordinating full-line Brownian
motion and then killing the jump process when it exits the half-line.
The latter gives a different nonlocal boundary problem; the odd
extension above does not prove its hard-wall limit. See
\href{https://arxiv.org/abs/1610.00872}{Kim--Song--Vondraček}
(introduction read) for this order-of-operations distinction. Nor is the
local Schrödinger current in the guide automatically the probability
current for this nonlocal Hamiltonian.

\hypertarget{consequence-for-state}{%
\subsection{6. Consequence for STATE}\label{consequence-for-state}}

For the positive-energy one-particle dispersion, bounded curvature
closes the quantum packet Cauchy argument with constants uniform in
\(c\). The fixed-\(\hbar\) form limit and the joint observable limit now
connect relativistic and nonrelativistic mechanics in the same
framework, with a separate spectral hard-wall example. The action unit
is supplied, and no new Yang--Mills or positive-action-necessity
conclusion follows.
\end{document}
